% Options for packages loaded elsewhere
\PassOptionsToPackage{unicode}{hyperref}
\PassOptionsToPackage{hyphens}{url}
\documentclass[
]{article}
\usepackage{ma2003b-notes}
\hypersetup{
  pdftitle={Principal Component Analysis - Lecture Notes},
  pdfauthor={MA2003B},
  hidelinks,
  pdfcreator={LaTeX via pandoc}}

\title{Principal Component Analysis - Lecture Notes}
\author{MA2003B}
\date{}

\begin{document}
\maketitle

\textbf{Principal Component Analysis (PCA)}

Dimensionality Reduction, Spectral Decomposition, and Geometry

MA2003B --- Application of Multivariate Methods in Data Science

Tecnológico de Monterrey

\begin{center}\rule{0.5\linewidth}{0.5pt}\end{center}

\textbf{ }

\section{Course and Chapter Overview}

\textbf{ }

\subsection{Course Context}

Principal Component Analysis (PCA) is the primary unsupervised
dimensionality reduction technique in multivariate data science.
Introduced by Karl Pearson (1901) and generalized by Harold Hotelling
(1933), PCA transforms a system of \(p\) correlated variables into
\(k \ll p\) mutually orthogonal components that capture the maximum
possible variance.

\textbf{ }

\subsection{Learning Objectives}

By the end of this chapter, students will be able to:

\begin{itemize}
\item
  Identify appropriate practical and scientific use cases for PCA.
\item
  Understand the geometric interpretation of PCA as an orthogonal
  rotation along the principal axes of the covariance ellipsoid.
\item
  Derive the principal component formulation via spectral decomposition
  of \(\mathbf{S}\) and \(\mathbf{R}\) using Lagrange multipliers.
\item
  Calculate component loadings, individual scores, and proportions of
  variance explained.
\item
  Apply rigorous component retention criteria (Kaiser criterion, Scree
  plots, Horn's Parallel Analysis).
\item
  Construct and interpret 2D/3D PCA Biplots.
\item
  Implement end-to-end PCA pipelines in Python using
  \texttt{scikit-learn}.
\end{itemize}

---

\textbf{ }

\section{3.1 Cases Where PCA is Used}

\textbf{ }

\subsection{Primary Applications}

\begin{enumerate}
\item
  \textbf{}Multicollinearity Elimination:\textbf{} Replaces highly
  correlated regression predictors with uncorrelated principal
  components (Principal Component Regression, PCR).
\item
  \textbf{}High-Dimensional Data Compression:\textbf{} Reduces memory
  footprint and computational cost in downstream machine learning
  algorithms while retaining \(> 80\%\) of systemic variance.
\item
  \textbf{}Exploratory 2D/3D Visualization:\textbf{} Projects
  high-dimensional datasets (\(p \geq 10\)) onto principal score planes
  to identify natural clustering, trajectories, and outliers.
\item
  \textbf{}Macroeconomic \& Financial Risk Factor Modeling:\textbf{}
  Extracts systematic market drivers (e.g. Yield curve level, slope,
  curvature) from multi-asset return streams.
\end{enumerate}

\textbf{Info: } \textbf{PCA vs. Factor Analysis:} PCA is a
variance-focused mathematical transformation of all observed variance
(common + unique). In contrast, Factor Analysis posits an underlying
measurement model focusing strictly on shared common variance
(\(\mathbf{\Sigma} = \mathbf{\Lambda}\mathbf{\Lambda}^{T} + \mathbf{\Psi}\)).

---

\textbf{ }

\section{3.2 Geometrical Description of Major Components}

\textbf{ }

\subsection{Geometric Rotation of Axes}

Let \(\mathbf{x}_{1},\ldots,\mathbf{x}_{n} \in \mathbb{R}^{p}\) be a
sample of observations centered at
\(\overline{\mathbf{x}} = \mathbf{0}\).

\begin{itemize}
\item
  The \textbf{first principal component}
  \(Y_{1} = \mathbf{a}_{1}^{T}\mathbf{X}\) defines the direction in
  \(\mathbb{R}^{p}\) along which the orthogonal projection of the data
  points exhibits \textbf{maximum dispersion (variance)}.
\item
  The \textbf{second principal component}
  \(Y_{2} = \mathbf{a}_{2}^{T}\mathbf{X}\) defines the direction of
  maximum remaining variance, subject to being \textbf{orthogonal} to
  \(\mathbf{a}_{1}\) (\(\mathbf{a}_{2}^{T}\mathbf{a}_{1} = 0\)).
\item
  Geometrically, the eigenvectors of \(\mathbf{S}\) specify the
  orientation of the principal axes of the constant-density ellipsoid
  \(\{\mathbf{x}:\mathbf{x}^{T}\mathbf{S}^{- 1}\mathbf{x} = c^{2}\}\).
\end{itemize}

---

\textbf{ }

\section{3.3 Mathematical Formulation \& Spectral Decomposition}

\textbf{ }

\subsection{Derivation via Lagrange Multipliers}

Let \(\mathbf{X}\) be a random vector with covariance matrix
\(\mathbf{\Sigma}\) (or sample covariance \(\mathbf{S}\)). We seek a
linear combination \(Y_{1} = \mathbf{a}_{1}^{T}\mathbf{X}\) that
maximizes:

\[\operatorname{Var}(Y_{1}) = \mathbf{a}_{1}^{T}\mathbf{\Sigma}\mathbf{a}_{1}\quad\text{ subject to  }\mathbf{a}_{1}^{T}\mathbf{a}_{1} = 1\]

The Lagrangian is:
\[\mathcal{L}(\mathbf{a}_{1},\lambda_{1}) = \mathbf{a}_{1}^{T}\mathbf{\Sigma}\mathbf{a}_{1} - \lambda_{1}\left( \mathbf{a}_{1}^{T}\mathbf{a}_{1} - 1 \right)\]

Setting the gradient to zero:
\[\frac{\partial\mathcal{L}}{\partial\mathbf{a}_{1}} = 2\mathbf{\Sigma}\mathbf{a}_{1} - 2\lambda_{1}\mathbf{a}_{1} = \mathbf{0} \Longrightarrow \mathbf{\Sigma}\mathbf{a}_{1} = \lambda_{1}\mathbf{a}_{1}\]

Thus:

\begin{itemize}
\item
  \(\lambda_{1}\) is the \textbf{largest eigenvalue} of
  \(\mathbf{\Sigma}\).
\item
  \(\mathbf{a}_{1}\) is the corresponding \textbf{normalized
  eigenvector} (\(|\mathbf{a}_{1}| = 1\)).
\item
  Maximum variance equals the eigenvalue:
  \(\operatorname{Var}(Y_{1}) = \mathbf{a}_{1}^{T}\mathbf{\Sigma}\mathbf{a}_{1} = \lambda_{1}\).
\end{itemize}

For the \(j\)-th component:
\[\mathbf{\Sigma}\mathbf{a}_{j} = \lambda_{j}\mathbf{a}_{j},\quad\lambda_{1} \geq \lambda_{2} \geq \ldots \geq \lambda_{p} \geq 0\]

\textbf{ }

\subsection{\texorpdfstring{Spectral Decomposition of Covariance
(\(\mathbf{S}\)) and Correlation
(\(\mathbf{R}\))}{Spectral Decomposition of Covariance (\textbackslash mathbf\{S\}) and Correlation (\textbackslash mathbf\{R\})}}

\[\mathbf{S} = \mathbf{V}\mathbf{\Lambda}\mathbf{V}^{T} = \sum_{j = 1}^{p}\lambda_{j}\mathbf{v}_{j}\mathbf{v}_{j}^{T}\]

where
\(\mathbf{V} = \left\lbrack \mathbf{v}_{1},\ldots,\mathbf{v}_{p} \right\rbrack\)
is the orthogonal eigenvector matrix
(\(\mathbf{V}^{T}\mathbf{V} = \mathbf{I}\)), and
\(\mathbf{\Lambda} = \operatorname{diag}(\lambda_{1},\ldots,\lambda_{p})\).

\textbf{ }

\subsection{Total Variance and Variance Explained}

\[\text{ Total Variance } = \operatorname{tr}(\mathbf{S}) = \sum_{j = 1}^{p}s_{jj} = \sum_{j = 1}^{p}\lambda_{j}\]

The proportion of total sample variance explained by the \(j\)-th
principal component is:
\[\text{ Variance Proportion}_{j} = \frac{\lambda_{j}}{\sum_{k = 1}^{p}\lambda_{k}}\]

\textbf{ }

\subsection{Component Loadings (Correlations)}

The correlation between the original variable \(X_{k}\) and the \(j\)-th
component \(Y_{j}\) is:
\[\operatorname{Corr}(X_{k},Y_{j}) = \frac{a_{jk}\sqrt{\lambda_{j}}}{\sqrt{s_{kk}}}\]

When PCA is performed on the \textbf{correlation matrix} \(\mathbf{R}\)
(where \(s_{kk} = 1\)):
\[\operatorname{Corr}(X_{k},Y_{j}) = a_{jk}\sqrt{\lambda_{j}}\]

---

\textbf{ }

\section{3.4 Determination of the Number of Major Components}

Several complementary rules guide component retention:

{\def\LTcaptype{none} % do not increment counter
\begin{longtable}[]{@{}
  >{\raggedright\arraybackslash}p{(\linewidth - 4\tabcolsep) * \real{0.2858}}
  >{\raggedright\arraybackslash}p{(\linewidth - 4\tabcolsep) * \real{0.4286}}
  >{\raggedright\arraybackslash}p{(\linewidth - 4\tabcolsep) * \real{0.2858}}@{}}
\toprule\noalign{}
\begin{minipage}[b]{\linewidth}\raggedright
\textbf{Criterion}
\end{minipage} & \begin{minipage}[b]{\linewidth}\raggedright
\textbf{Mathematical Rule}
\end{minipage} & \begin{minipage}[b]{\linewidth}\raggedright
\textbf{Best Use Context}
\end{minipage} \\
\midrule\noalign{}
\endhead
\bottomrule\noalign{}
\endlastfoot
Kaiser-Guttman Rule & Retain components where \(\lambda_{j} > 1.0\) (for
\(\mathbf{R}\)) or \(\lambda_{j} > \overline{\lambda}\) (for
\(\mathbf{S}\)) & Standard correlation matrix PCA \\
Cattell Scree Plot & Visually locate the ``elbow'' (inflection point)
where eigenvalues level off & Exploratory visual inspection \\
Cumulative Variance & Retain minimum \(k\) such that
\(\sum_{j = 1}^{k}\frac{\lambda_{j}}{\sum}\lambda_{m} \geq 70\% - 85\%\)
& Data compression goals \\
Horn's Parallel Analysis & Retain \(\lambda_{j}\) that exceeds the 95th
percentile of eigenvalues from random uncorrelated noise & Rigorous
statistical benchmark \\
\end{longtable}
}

---

\textbf{ }

\section{3.5 Biplots and Interpretation in Python}

\textbf{ }

\subsection{The PCA Biplot (Gabriel 1971)}

A Biplot simultaneously displays two sets of coordinates on the first
two principal axes:

\begin{enumerate}
\item
  \textbf{}Observation Scores (\(Y_{i1},Y_{i2}\)):\textbf{} Plotted as
  scatter points showing sample clustering.
\item
  \textbf{}Variable Loading Vectors
  (\(\mathbf{a}_{1},\mathbf{a}_{2}\)):\textbf{} Plotted as vectors
  emanating from the origin.

  \begin{itemize}
  \item
    Vector length indicates the strength of representation in the 2D
    subspace.
  \item
    The cosine of the angle between two variable vectors approximates
    their correlation.
  \end{itemize}
\end{enumerate}

---

\textbf{ }

\section{Practical Python Implementation}

\begin{Shaded}
\begin{Highlighting}[]
\ImportTok{import}\NormalTok{ numpy }\ImportTok{as}\NormalTok{ np}
\ImportTok{import}\NormalTok{ pandas }\ImportTok{as}\NormalTok{ pd}
\ImportTok{import}\NormalTok{ matplotlib.pyplot }\ImportTok{as}\NormalTok{ plt}
\ImportTok{from}\NormalTok{ sklearn.preprocessing }\ImportTok{import}\NormalTok{ StandardScaler}
\ImportTok{from}\NormalTok{ sklearn.decomposition }\ImportTok{import}\NormalTok{ PCA}

\CommentTok{\# 1. Standardize Features}
\NormalTok{scaler }\OperatorTok{=}\NormalTok{ StandardScaler()}
\NormalTok{X\_scaled }\OperatorTok{=}\NormalTok{ scaler.fit\_transform(df.iloc[:, }\DecValTok{1}\NormalTok{:])}

\CommentTok{\# 2. Fit Full PCA}
\NormalTok{pca }\OperatorTok{=}\NormalTok{ PCA()}
\NormalTok{scores }\OperatorTok{=}\NormalTok{ pca.fit\_transform(X\_scaled)}
\NormalTok{eigenvalues }\OperatorTok{=}\NormalTok{ pca.explained\_variance\_}
\NormalTok{variance\_ratio }\OperatorTok{=}\NormalTok{ pca.explained\_variance\_ratio\_}

\CommentTok{\# 3. Scree Plot \& Kaiser Criterion}
\NormalTok{plt.figure(figsize}\OperatorTok{=}\NormalTok{(}\DecValTok{8}\NormalTok{, }\DecValTok{4}\NormalTok{))}
\NormalTok{plt.plot(}\BuiltInTok{range}\NormalTok{(}\DecValTok{1}\NormalTok{, }\BuiltInTok{len}\NormalTok{(eigenvalues) }\OperatorTok{+} \DecValTok{1}\NormalTok{), eigenvalues, }\StringTok{\textquotesingle{}o{-}\textquotesingle{}}\NormalTok{, color}\OperatorTok{=}\StringTok{\textquotesingle{}navy\textquotesingle{}}\NormalTok{, lw}\OperatorTok{=}\DecValTok{2}\NormalTok{)}
\NormalTok{plt.axhline(}\FloatTok{1.0}\NormalTok{, color}\OperatorTok{=}\StringTok{\textquotesingle{}red\textquotesingle{}}\NormalTok{, ls}\OperatorTok{=}\StringTok{\textquotesingle{}{-}{-}\textquotesingle{}}\NormalTok{, label}\OperatorTok{=}\StringTok{\textquotesingle{}Kaiser Threshold (λ = 1.0)\textquotesingle{}}\NormalTok{)}
\NormalTok{plt.xlabel(}\StringTok{\textquotesingle{}Principal Component Index\textquotesingle{}}\NormalTok{)}
\NormalTok{plt.ylabel(}\StringTok{\textquotesingle{}Eigenvalue (Variance Explained)\textquotesingle{}}\NormalTok{)}
\NormalTok{plt.title(}\StringTok{\textquotesingle{}PCA Scree Plot\textquotesingle{}}\NormalTok{)}
\NormalTok{plt.legend()}
\NormalTok{plt.show()}

\CommentTok{\# 4. Component Loadings DataFrame}
\NormalTok{loadings }\OperatorTok{=}\NormalTok{ pd.DataFrame(}
\NormalTok{    pca.components\_[:}\DecValTok{3}\NormalTok{].T }\OperatorTok{*}\NormalTok{ np.sqrt(pca.explained\_variance\_[:}\DecValTok{3}\NormalTok{]),}
\NormalTok{    index}\OperatorTok{=}\NormalTok{df.columns[}\DecValTok{1}\NormalTok{:],}
\NormalTok{    columns}\OperatorTok{=}\NormalTok{[}\StringTok{\textquotesingle{}PC1\textquotesingle{}}\NormalTok{, }\StringTok{\textquotesingle{}PC2\textquotesingle{}}\NormalTok{, }\StringTok{\textquotesingle{}PC3\textquotesingle{}}\NormalTok{]}
\NormalTok{)}
\BuiltInTok{print}\NormalTok{(}\StringTok{"=== Component Loadings Matrix ==="}\NormalTok{)}
\BuiltInTok{print}\NormalTok{(loadings.}\BuiltInTok{round}\NormalTok{(}\DecValTok{3}\NormalTok{))}
\end{Highlighting}
\end{Shaded}


\end{document}



